\documentclass[11pt]{article}
\usepackage{amsmath,amssymb}
\title{Room 5 (P5 W\_T non-D-finite): moderator summary}
\date{2026-09-26}
\begin{document}
\maketitle

Two seats (Hilbert, Erd\H os), each attempting a route different from the already-tried
catalytic-variable kernel method (\texttt{P5\_v2\_note.tex}).

\textbf{Hilbert: NOTHING WORKED, three routes precisely blocked.} A contour/Cauchy-integral kernel
gives only a tautological restatement, since the reducibility relation (division by $2\mu_T$) has
an infinite, non-compact index set, unlike the finite-state travel block behind $U,V,G^T_0$. A
special-function (theta/$q$-Bessel) substitution fails: $\zeta_T$ is fixed, not a deforming/
confluent parameter, and the carry's length weight is a min-plus optimum, not a closed quadratic
form. A rationality transfer via the quasi-isometry to $\mathrm{DL}(5,5)$ fails because that
rationality result is for finite-lamp lamplighter groups, while $H_T$'s lamp group $A_T$ has rank 2
and is not finitely generated -- the QI does not preserve the exact generating function.

\textbf{Erd\H os: PARTIAL, same conclusion via different attempts.} A two-integer
(real/imaginary-coordinate) re-encoding of the catalytic step, via the monic recursion for
$\eta=c\zeta$, relocates rather than removes the obstruction: the integer-pair recursion itself
grows geometrically ($\sim c^h$) per step, reproducing the same unbounded-carry fact (independently
reconfirming, in a different basis, that $A_T$ is infinite -- consistent with, not beyond, the
already-known result). A tower of finite-radius truncations $F_T^{(N)}(x)$ (each individually
rational/D-finite, converging coefficientwise to $F_T$) is a genuinely new formal object -- paper1
only states the analogy to Akiyama--Frougny--Sakarovitch in words -- but provides no new leverage;
it repackages the unproved carry-width growth law rather than resolving it.

\textbf{Room verdict: NO PROGRESS,} but with two independently-reached, precisely-diagnosed
structural obstructions (rank-2 non-f.g.\ lamp group defeats the QI-rationality shortcut; the
min-plus/non-contracting-base structure defeats every continuation route tried) and one small,
honest formal artifact (the truncation tower) that clarifies the picture without advancing it. P5
stands exactly where \texttt{P5\_note.tex}/\texttt{P5\_v2\_note.tex} left it: OPEN, with Gaps 1--3
now attacked from two more angles, both closed off for stated reasons.

\end{document}
